% assets/w03-latlon-ned.svg — 위도·경도가 무엇이고, 그 자리에 NED 가 어떻게 서는가
%
%   bash _tools/tikz2svg.sh assets/src/w03-latlon-ned.tex assets/w03-latlon-ned.svg
%
% 요점
%   위도 phi 는 적도에서 북쪽으로 잰 각, 경도 lambda 는 본초자오선에서 동쪽으로 잰 각.
%   그 점에 지구를 **접하는 평면**을 세우면 N-E-D 가 된다. D 는 지구 중심을 향한다.
%
% 투영
%   정사영. 카메라가 경도 LC 위, 적도에서 T 만큼 올라간 자리에 있다.
%     screen_x = R cos(phi) sin(lambda - LC)
%     screen_y = R ( cos(T) sin(phi) - sin(T) cos(phi) cos(lambda - LC) )
%   LC = 0 이면 본초자오선이 화면 한가운데 직선이 되어 극축과 겹친다 (2026-09-18 고침, LC = -25)
%   3차원 단위벡터 (vx,vy,vz) 는
%     screen_x = vy,   screen_y = -sin(T) vx + cos(T) vz

\documentclass[border=6pt]{standalone}
% capstone-style.tex — 이 볼트의 TikZ 그림이 공유하는 스타일
%
% 쓰는 법:  % capstone-style.tex — 이 볼트의 TikZ 그림이 공유하는 스타일
%
% 쓰는 법:  % capstone-style.tex — 이 볼트의 TikZ 그림이 공유하는 스타일
%
% 쓰는 법:  \input{capstone-style}   (assets/src/ 안에서)
% 빌드:     bash _tools/tikz2svg.sh assets/src/w02-node-topic.tex assets/w02-node-topic.svg
%
% 왜 TikZ 인가
%   손으로 SVG 좌표를 쓰면 화살촉·선 굵기·글자 크기가 그림마다 미묘하게 달라진다.
%   그것이 "자료가 어설퍼 보인다"의 정체다. 스타일을 한 번만 정하고 상대좌표로
%   배치하면 좌표를 고쳐도 그림이 무너지지 않는다.
%   수식은 본문과 같은 TeX 글꼴로 나온다.
%
% 한글
%   ko.TeX(DVI 경로)를 쓴다. XeLaTeX 은 TikZ 그래픽을 PDF special 로 내보내는데
%   dvisvgm 이 그것을 받으려면 Ghostscript < 10.01 이 필요하다. 이 PC 에는
%   10.07 이 깔려 있어 그 경로가 막힌다. latex -> DVI -> dvisvgm 이 유일하게
%   동작하는 조합이다.

\usepackage[hangul]{kotex}
\usepackage{tikz}
\usepackage{amsmath}
\usepackage{xcolor}
\usetikzlibrary{arrows.meta, positioning, calc, fit, backgrounds, decorations.pathmorphing}

% ---- 색 --------------------------------------------------------------------
% 강조는 하나만. 나머지는 검정. 흑백 인쇄에서도 읽혀야 한다.
\definecolor{ink}{HTML}{1A1A1A}
\definecolor{soft}{HTML}{666666}
\definecolor{accent}{HTML}{7C3AED}
\definecolor{warn}{HTML}{B91C1C}

% 계층 색 — Simulink 모델의 블록 색과 같은 규칙을 쓴다
\definecolor{guid}{HTML}{1D4ED8}    % 유도
\definecolor{ctrl}{HTML}{EA580C}    % 제어
\definecolor{thr}{HTML}{D97706}     % 추진기
\definecolor{plant}{HTML}{16A34A}   % 운동모델
\definecolor{nav}{HTML}{7C3AED}     % 항법 · 신호처리
\definecolor{pathc}{HTML}{0D9488}   % 계획 경로 · 항적
\definecolor{errc}{HTML}{C2410C}    % 오차
\definecolor{flow}{HTML}{0369A1}    % 조류 · 바람

% ---- 선과 화살촉 -----------------------------------------------------------
\tikzset{
  sig/.style   = {ink, line width=0.5pt, -{Stealth[length=4pt, width=3pt]}},
  wire/.style  = {ink, line width=0.5pt},
  fb/.style    = {ink, line width=0.5pt, densely dashed,
                  -{Stealth[length=4pt, width=3pt]}},
  asig/.style  = {accent, line width=0.5pt, densely dashed,
                  -{Stealth[length=4pt, width=3pt]}},
  % 블록 — 흰 바탕, 직각, 검은 테두리. 색은 테두리에만 준다
  blk/.style   = {draw=ink, line width=0.55pt, fill=white,
                  minimum width=18mm, minimum height=8.5mm, inner sep=2pt, font=\small},
  gblk/.style  = {blk, draw=guid},
  cblk/.style  = {blk, draw=ctrl},
  tblk/.style  = {blk, draw=thr},
  pblk/.style  = {blk, draw=plant},
  nblk/.style  = {blk, draw=nav},
  % 합산점 · 분기점
  sum/.style   = {draw=ink, line width=0.55pt, fill=white, circle,
                  inner sep=0pt, minimum size=4.6mm, font=\scriptsize},
  dot/.style   = {ink, fill=ink, circle, inner sep=0pt, minimum size=1.5mm},
  % 글자
  sname/.style = {font=\small, inner sep=1.5pt},
  note/.style  = {font=\scriptsize, soft, inner sep=1.5pt},
  wnote/.style = {font=\scriptsize, warn, inner sep=1.5pt},
  anote/.style = {font=\scriptsize, accent, inner sep=1.5pt},
  axis/.style  = {ink, line width=0.5pt},
  tick/.style  = {font=\scriptsize, soft},
  % 묶음 상자 — 서브시스템 표시
  grp/.style   = {draw=soft, line width=0.4pt, densely dotted, rounded corners=1mm,
                  inner sep=3mm},
}
   (assets/src/ 안에서)
% 빌드:     bash _tools/tikz2svg.sh assets/src/w02-node-topic.tex assets/w02-node-topic.svg
%
% 왜 TikZ 인가
%   손으로 SVG 좌표를 쓰면 화살촉·선 굵기·글자 크기가 그림마다 미묘하게 달라진다.
%   그것이 "자료가 어설퍼 보인다"의 정체다. 스타일을 한 번만 정하고 상대좌표로
%   배치하면 좌표를 고쳐도 그림이 무너지지 않는다.
%   수식은 본문과 같은 TeX 글꼴로 나온다.
%
% 한글
%   ko.TeX(DVI 경로)를 쓴다. XeLaTeX 은 TikZ 그래픽을 PDF special 로 내보내는데
%   dvisvgm 이 그것을 받으려면 Ghostscript < 10.01 이 필요하다. 이 PC 에는
%   10.07 이 깔려 있어 그 경로가 막힌다. latex -> DVI -> dvisvgm 이 유일하게
%   동작하는 조합이다.

\usepackage[hangul]{kotex}
\usepackage{tikz}
\usepackage{amsmath}
\usepackage{xcolor}
\usetikzlibrary{arrows.meta, positioning, calc, fit, backgrounds, decorations.pathmorphing}

% ---- 색 --------------------------------------------------------------------
% 강조는 하나만. 나머지는 검정. 흑백 인쇄에서도 읽혀야 한다.
\definecolor{ink}{HTML}{1A1A1A}
\definecolor{soft}{HTML}{666666}
\definecolor{accent}{HTML}{7C3AED}
\definecolor{warn}{HTML}{B91C1C}

% 계층 색 — Simulink 모델의 블록 색과 같은 규칙을 쓴다
\definecolor{guid}{HTML}{1D4ED8}    % 유도
\definecolor{ctrl}{HTML}{EA580C}    % 제어
\definecolor{thr}{HTML}{D97706}     % 추진기
\definecolor{plant}{HTML}{16A34A}   % 운동모델
\definecolor{nav}{HTML}{7C3AED}     % 항법 · 신호처리
\definecolor{pathc}{HTML}{0D9488}   % 계획 경로 · 항적
\definecolor{errc}{HTML}{C2410C}    % 오차
\definecolor{flow}{HTML}{0369A1}    % 조류 · 바람

% ---- 선과 화살촉 -----------------------------------------------------------
\tikzset{
  sig/.style   = {ink, line width=0.5pt, -{Stealth[length=4pt, width=3pt]}},
  wire/.style  = {ink, line width=0.5pt},
  fb/.style    = {ink, line width=0.5pt, densely dashed,
                  -{Stealth[length=4pt, width=3pt]}},
  asig/.style  = {accent, line width=0.5pt, densely dashed,
                  -{Stealth[length=4pt, width=3pt]}},
  % 블록 — 흰 바탕, 직각, 검은 테두리. 색은 테두리에만 준다
  blk/.style   = {draw=ink, line width=0.55pt, fill=white,
                  minimum width=18mm, minimum height=8.5mm, inner sep=2pt, font=\small},
  gblk/.style  = {blk, draw=guid},
  cblk/.style  = {blk, draw=ctrl},
  tblk/.style  = {blk, draw=thr},
  pblk/.style  = {blk, draw=plant},
  nblk/.style  = {blk, draw=nav},
  % 합산점 · 분기점
  sum/.style   = {draw=ink, line width=0.55pt, fill=white, circle,
                  inner sep=0pt, minimum size=4.6mm, font=\scriptsize},
  dot/.style   = {ink, fill=ink, circle, inner sep=0pt, minimum size=1.5mm},
  % 글자
  sname/.style = {font=\small, inner sep=1.5pt},
  note/.style  = {font=\scriptsize, soft, inner sep=1.5pt},
  wnote/.style = {font=\scriptsize, warn, inner sep=1.5pt},
  anote/.style = {font=\scriptsize, accent, inner sep=1.5pt},
  axis/.style  = {ink, line width=0.5pt},
  tick/.style  = {font=\scriptsize, soft},
  % 묶음 상자 — 서브시스템 표시
  grp/.style   = {draw=soft, line width=0.4pt, densely dotted, rounded corners=1mm,
                  inner sep=3mm},
}
   (assets/src/ 안에서)
% 빌드:     bash _tools/tikz2svg.sh assets/src/w02-node-topic.tex assets/w02-node-topic.svg
%
% 왜 TikZ 인가
%   손으로 SVG 좌표를 쓰면 화살촉·선 굵기·글자 크기가 그림마다 미묘하게 달라진다.
%   그것이 "자료가 어설퍼 보인다"의 정체다. 스타일을 한 번만 정하고 상대좌표로
%   배치하면 좌표를 고쳐도 그림이 무너지지 않는다.
%   수식은 본문과 같은 TeX 글꼴로 나온다.
%
% 한글
%   ko.TeX(DVI 경로)를 쓴다. XeLaTeX 은 TikZ 그래픽을 PDF special 로 내보내는데
%   dvisvgm 이 그것을 받으려면 Ghostscript < 10.01 이 필요하다. 이 PC 에는
%   10.07 이 깔려 있어 그 경로가 막힌다. latex -> DVI -> dvisvgm 이 유일하게
%   동작하는 조합이다.

\usepackage[hangul]{kotex}
\usepackage{tikz}
\usepackage{amsmath}
\usepackage{xcolor}
\usetikzlibrary{arrows.meta, positioning, calc, fit, backgrounds, decorations.pathmorphing}

% ---- 색 --------------------------------------------------------------------
% 강조는 하나만. 나머지는 검정. 흑백 인쇄에서도 읽혀야 한다.
\definecolor{ink}{HTML}{1A1A1A}
\definecolor{soft}{HTML}{666666}
\definecolor{accent}{HTML}{7C3AED}
\definecolor{warn}{HTML}{B91C1C}

% 계층 색 — Simulink 모델의 블록 색과 같은 규칙을 쓴다
\definecolor{guid}{HTML}{1D4ED8}    % 유도
\definecolor{ctrl}{HTML}{EA580C}    % 제어
\definecolor{thr}{HTML}{D97706}     % 추진기
\definecolor{plant}{HTML}{16A34A}   % 운동모델
\definecolor{nav}{HTML}{7C3AED}     % 항법 · 신호처리
\definecolor{pathc}{HTML}{0D9488}   % 계획 경로 · 항적
\definecolor{errc}{HTML}{C2410C}    % 오차
\definecolor{flow}{HTML}{0369A1}    % 조류 · 바람

% ---- 선과 화살촉 -----------------------------------------------------------
\tikzset{
  sig/.style   = {ink, line width=0.5pt, -{Stealth[length=4pt, width=3pt]}},
  wire/.style  = {ink, line width=0.5pt},
  fb/.style    = {ink, line width=0.5pt, densely dashed,
                  -{Stealth[length=4pt, width=3pt]}},
  asig/.style  = {accent, line width=0.5pt, densely dashed,
                  -{Stealth[length=4pt, width=3pt]}},
  % 블록 — 흰 바탕, 직각, 검은 테두리. 색은 테두리에만 준다
  blk/.style   = {draw=ink, line width=0.55pt, fill=white,
                  minimum width=18mm, minimum height=8.5mm, inner sep=2pt, font=\small},
  gblk/.style  = {blk, draw=guid},
  cblk/.style  = {blk, draw=ctrl},
  tblk/.style  = {blk, draw=thr},
  pblk/.style  = {blk, draw=plant},
  nblk/.style  = {blk, draw=nav},
  % 합산점 · 분기점
  sum/.style   = {draw=ink, line width=0.55pt, fill=white, circle,
                  inner sep=0pt, minimum size=4.6mm, font=\scriptsize},
  dot/.style   = {ink, fill=ink, circle, inner sep=0pt, minimum size=1.5mm},
  % 글자
  sname/.style = {font=\small, inner sep=1.5pt},
  note/.style  = {font=\scriptsize, soft, inner sep=1.5pt},
  wnote/.style = {font=\scriptsize, warn, inner sep=1.5pt},
  anote/.style = {font=\scriptsize, accent, inner sep=1.5pt},
  axis/.style  = {ink, line width=0.5pt},
  tick/.style  = {font=\scriptsize, soft},
  % 묶음 상자 — 서브시스템 표시
  grp/.style   = {draw=soft, line width=0.4pt, densely dotted, rounded corners=1mm,
                  inner sep=3mm},
}


\begin{document}
\begin{tikzpicture}[x=1mm, y=1mm]

\pgfmathsetmacro{\R}{28}        % 지구 반지름 [mm]
\pgfmathsetmacro{\T}{22}        % 카메라 올라간 각 [deg]
\pgfmathsetmacro{\pP}{34}       % 점 P 의 위도 [deg]
\pgfmathsetmacro{\lP}{48}       % 점 P 의 경도 [deg]
\pgfmathsetmacro{\LC}{-25}      % 카메라 경도 [deg]

% 화면 좌표 — #1 = 위도, #2 = 경도
\newcommand{\SX}[2]{\R*cos(#1)*sin((#2)-\LC)}
\newcommand{\SY}[2]{\R*(cos(\T)*sin(#1) - sin(\T)*cos(#1)*cos((#2)-\LC))}

% =========================== 지구 ==========================================
\begin{scope}[shift={(0,0)}]

  \fill[white] (0,0) circle (\R);
  \draw[soft, line width=0.45pt] (0,0) circle (\R);

  % ---- 적도 : 앞(실선) · 뒤(점선) ----
  \draw[soft, line width=0.35pt, densely dotted, samples=80, domain=\LC+90:\LC+270, variable=\t]
       plot ({\SX{0}{\t}}, {\SY{0}{\t}});
  \draw[errc, line width=0.7pt, samples=80, domain=\LC-90:\LC+90, variable=\t]
       plot ({\SX{0}{\t}}, {\SY{0}{\t}});
  \node[font=\scriptsize, errc, anchor=north] at ({\SX{0}{-62}}, {\SY{0}{-62}-1.5}) {적도};

  % ---- 본초자오선 (경도 0) ----
  \draw[guid, line width=0.7pt, samples=60, domain=-90:90, variable=\t]
       plot ({\SX{\t}{0}}, {\SY{\t}{0}});
  \node[font=\scriptsize, guid, anchor=west] at ({\SX{-40}{0}+1.5}, {\SY{-40}{0}}) {본초자오선};

  % ---- P 의 자오선과 위도선 ----
  \draw[soft, line width=0.45pt, samples=60, domain=-90:90, variable=\t]
       plot ({\SX{\t}{\lP}}, {\SY{\t}{\lP}});
  \draw[soft, line width=0.35pt, densely dashed, samples=70, domain=\LC-90:\LC+90, variable=\t]
       plot ({\SX{\pP}{\t}}, {\SY{\pP}{\t}});

  % ---- 극축 ----
  \draw[soft, line width=0.35pt, densely dashed]
       ({\SX{-90}{0}}, {\SY{-90}{0}}) -- ({\SX{90}{0}}, {\SY{90}{0}});
  \node[font=\scriptsize, soft, anchor=south] at ({\SX{90}{0}}, {\SY{90}{0}+1}) {북극};

  % ---- 경도 각 : 적도를 따라 0 에서 lP 까지 ----
  \draw[accent, line width=0.9pt, samples=40, domain=0:\lP, variable=\t]
       plot ({0.52*\SX{0}{\t}}, {0.52*\SY{0}{\t}});
  \node[font=\small, accent] at ({0.66*\SX{0}{26}}, {0.66*\SY{0}{26}-1}) {$\lambda$};

  % ---- 위도 각 : P 의 자오선을 따라 적도에서 P 까지 ----
  \draw[ctrl, line width=0.9pt, samples=40, domain=0:\pP, variable=\t]
       plot ({0.52*\SX{\t}{\lP}}, {0.52*\SY{\t}{\lP}});
  \node[font=\small, ctrl] at ({0.70*\SX{18}{\lP}+1.5}, {0.70*\SY{18}{\lP}}) {$\varphi$};

  % ---- 중심에서 P 로, 그리고 적도 발까지 ----
  \draw[soft, line width=0.35pt, densely dashed]
       (0,0) -- ({\SX{\pP}{\lP}}, {\SY{\pP}{\lP}});
  \draw[soft, line width=0.35pt, densely dashed]
       (0,0) -- ({\SX{0}{\lP}}, {\SY{0}{\lP}});
  \fill[ink] (0,0) circle (0.7);

  % ---- 점 P ----
  \fill[ink] ({\SX{\pP}{\lP}}, {\SY{\pP}{\lP}}) circle (1.1);
  \node[font=\small, ink, anchor=south west]
       at ({\SX{\pP}{\lP}+1.2}, {\SY{\pP}{\lP}+0.8}) {$P$};

  \node[font=\small\bfseries, anchor=west] at (-\R,\R+12) {위도 $\varphi$ 와 경도 $\lambda$};
  \node[note, anchor=west, text width=62mm] at (-\R,\R+6)
       {\textcolor{ctrl}{\textbf{위도}} 는 적도에서 북쪽으로 잰 각,
        \textcolor{accent}{\textbf{경도}} 는 본초자오선에서 동쪽으로 잰 각};
\end{scope}

% =========================== 접평면과 NED ==================================
\begin{scope}[shift={(96,0)}]

  \pgfmathsetmacro{\Rr}{22}
  % P 의 3차원 단위벡터
  \pgfmathsetmacro{\px}{cos(\pP)*cos(\lP-\LC)}
  \pgfmathsetmacro{\py}{cos(\pP)*sin(\lP-\LC)}
  \pgfmathsetmacro{\pz}{sin(\pP)}
  % 북 : d(위치)/d(위도)
  \pgfmathsetmacro{\nx}{-sin(\pP)*cos(\lP-\LC)}
  \pgfmathsetmacro{\ny}{-sin(\pP)*sin(\lP-\LC)}
  \pgfmathsetmacro{\nz}{cos(\pP)}
  % 동 : d(위치)/d(경도)
  \pgfmathsetmacro{\ex}{-sin(\lP-\LC)}
  \pgfmathsetmacro{\ey}{cos(\lP-\LC)}
  \pgfmathsetmacro{\ez}{0}
  % 화면 투영
  \pgfmathsetmacro{\NX}{\ny}                       \pgfmathsetmacro{\NY}{-sin(\T)*\nx + cos(\T)*\nz}
  \pgfmathsetmacro{\EX}{\ey}                       \pgfmathsetmacro{\EY}{-sin(\T)*\ex + cos(\T)*\ez}
  \pgfmathsetmacro{\DX}{-\py}                      \pgfmathsetmacro{\DY}{-(-sin(\T)*\px + cos(\T)*\pz)}

  % 접평면 — 북·동 두 방향으로 펼친 사각형
  \pgfmathsetmacro{\HW}{20}
  \fill[guid!8]
       ({ \HW*\NX + \HW*\EX}, { \HW*\NY + \HW*\EY}) --
       ({ \HW*\NX - \HW*\EX}, { \HW*\NY - \HW*\EY}) --
       ({-\HW*\NX - \HW*\EX}, {-\HW*\NY - \HW*\EY}) --
       ({-\HW*\NX + \HW*\EX}, {-\HW*\NY + \HW*\EY}) -- cycle;
  \draw[soft, line width=0.4pt]
       ({ \HW*\NX + \HW*\EX}, { \HW*\NY + \HW*\EY}) --
       ({ \HW*\NX - \HW*\EX}, { \HW*\NY - \HW*\EY}) --
       ({-\HW*\NX - \HW*\EX}, {-\HW*\NY - \HW*\EY}) --
       ({-\HW*\NX + \HW*\EX}, {-\HW*\NY + \HW*\EY}) -- cycle;
  \node[note, anchor=west] at ({-\HW*\NX + \HW*\EX + 1}, {-\HW*\NY + \HW*\EY - 3})
       {$P$ 에서 지구에 접하는 평면};

  % NED 축
  \draw[errc, line width=1.0pt, -{Stealth[length=4.5pt,width=3.5pt]}]
       (0,0) -- ({\Rr*\NX}, {\Rr*\NY});
  \node[font=\small, errc] at ({(\Rr+5)*\NX}, {(\Rr+5)*\NY}) {$N$};
  \draw[errc, line width=1.0pt, -{Stealth[length=4.5pt,width=3.5pt]}]
       (0,0) -- ({\Rr*\EX}, {\Rr*\EY});
  \node[font=\small, errc] at ({(\Rr+5)*\EX}, {(\Rr+5)*\EY}) {$E$};
  \draw[errc, line width=1.0pt, -{Stealth[length=4.5pt,width=3.5pt]}]
       (0,0) -- ({0.8*\Rr*\DX}, {0.8*\Rr*\DY});
  \node[font=\small, errc] at ({(0.8*\Rr+5)*\DX}, {(0.8*\Rr+5)*\DY}) {$D$};

  \fill[ink] (0,0) circle (1.1);
  \node[font=\small, ink, anchor=south east] at (-1.2,1.0) {$P$};

  \node[font=\small\bfseries, anchor=west] at (-32,\R+12) {그 자리에 세우는 NED};
  \node[note, anchor=west, text width=62mm] at (-32,\R+6)
       {$N$ 은 자오선을 따라 북쪽, $E$ 는 위도선을 따라 동쪽,
        $D$ 는 \textbf{지구 중심}을 향한다};

  \node[note, anchor=west, text width=64mm] at (-32,-\R-4)
       {배가 수 km 안에서만 움직이면 이 평면을 평평하다고 보아도 된다.
        그것이 flat-earth 근사이며, 위경도 차이에 곡률반경을 곱해 미터로 바꾼다};
\end{scope}

% =========================== 변환식 ========================================
\node[font=\small, anchor=west] at (-28,-\R-22)
     {$\begin{aligned}
        \Delta N &= (\varphi - \varphi_0)\,R_M, \qquad
        \Delta E  = (\lambda - \lambda_0)\,R_N\cos\varphi_0, \qquad
        \Delta D  = -(h - h_0)
       \end{aligned}$};

\node[note, anchor=west, text width=150mm] at (-28,-\R-32)
     {$R_M$ 자오선 곡률반경, $R_N$ 묘유선 곡률반경. 경도 쪽에 $\cos\varphi_0$ 가 붙는 이유는
      위도가 높아질수록 경도 1도의 실제 거리가 짧아지기 때문이다 — 극에서는 0 이 된다};

\end{tikzpicture}
\end{document}
